% Generated draft: requires mathematical and source review.
\documentclass[11pt]{article}
\usepackage[margin=25mm]{geometry}
\usepackage{fontspec}
\setmainfont{Latin Modern Roman}
\usepackage{amsmath,amssymb,mathtools,unicode-math}
\setmathfont{Latin Modern Math}
\usepackage{enumitem,needspace,xurl,hyperref}
\setlength{\parindent}{0pt}
\setlength{\parskip}{5pt}
\setlength{\emergencystretch}{3em}
\newcommand{\N}{\mathbb{N}}
\newcommand{\Z}{\mathbb{Z}}
\newcommand{\Q}{\mathbb{Q}}
\newcommand{\R}{\mathbb{R}}
\newcommand{\C}{\mathbb{C}}
\newcommand{\F}{\mathbb{F}}
\newcommand{\A}{\mathbb{A}}
\newcommand{\PP}{\mathbb{P}}
\newcommand{\E}{\mathbb{E}}
\newcommand{\eps}{\varepsilon}
\newcommand{\dd}{\,\mathrm{d}}
\begin{document}
\section*{UnsolvedMath ID: 30004108}
\textbf{Source problem code:} catalogue entry
\subsection{Definitions and mathematical statement}
Let \(G\) be a finite group and let \(p\) be a prime. Under the usual \(p\)-adic interpretation of the source notation, \(\Z_p\) denotes the ring of \(p\)-adic integers. Let \(b\) be a block idempotent of \(\Z_pG\), that is, a primitive central idempotent, and assume that \(b\) has positive defect, meaning that its defect groups are nontrivial.

For a strict nonempty chain of nontrivial \(p\)-subgroups
\[
\sigma=(P_0<P_1<\cdots<P_r),\qquad 1\ne P_i\le G,
\]
conjugation by \(G\) acts on its terms. Its stabilizer is
\[
G_\sigma=\bigcap_{i=0}^{r}N_G(P_i),\qquad
N_G(P)=\{g\in G:gPg^{-1}=P\}.
\]
These definitions describe the usual strict-chain interpretation; the supplied record does not specify its chain convention or whether the empty chain is included. Write \([\mathcal P(G)/G]\) formally for a choice of representatives of the conjugacy classes of chains intended by the source. For a ring \(A\), \(\operatorname{Mat}_n(A)\) denotes the ring of \(n\times n\) matrices over \(A\), and \(|G:G_\sigma|\) is the subgroup index.

\textbf{Recorded cancellation question.} For every such \(G,p,b\), the supplied record asks whether
\[
\sum_{\sigma\in[\mathcal P(G)/G]}
(-1)^{|\sigma|}
\operatorname{Mat}_{|G:G_\sigma|}
\bigl(\Q G_\sigma b_\sigma\bigr)=0.
\]
\textbf{Unresolved qualifications.} This is a transcription of the intended target, not yet a well-defined algebraic equality. The record does not define \(b_\sigma\), specify the additive group or other structure in which the matrix-algebra terms are summed, or define the zero on the right. It also does not specify whether \(|\sigma|\) counts subgroup terms or strict inclusions, or whether an empty chain contributes. Finally, the formula explicitly uses \(\Q\), whereas \(b\) lies in a \(p\)-adic group ring and the supplied background describes a \(p\)-adic assertion. No construction placing \(b_\sigma\) in \(\Q G_\sigma\) is supplied. Replacing \(\Q\) by \(\Q_p\), choosing a Grothendieck group, or assigning local block idempotents would require additional source evidence. These choices remain unresolved.
\subsection{Short English statement}
The recorded question asks whether, for every finite group, prime, and positive-defect block, an alternating sum of matrix algebras attached to conjugacy classes of chains of nontrivial \(p\)-subgroups vanishes. The local block terms, coefficient field, chain conventions, and meaning of the sum require clarification before this becomes a precise conjecture. The supplied status is \emph{partially solved}; this draft does not independently verify that assessment.
\subsection{Sources}
\begin{itemize}
\item Local UnsolvedMath registry, permanent ID \(30004108\), title \emph{Blockwise Chain Cancellation for Finite Groups}, exact source problem code \texttt{catalogue entry}. The supplied record gives proposed year 2019 and names no proposer.
\item Supplied report citation: \emph{Representations of Finite Groups} (2020), Oberwolfach Reports, pp.\ 841--895, DOI \texttt{10.4171/owr/2019/14}. Source URL: \url{https://doi.org/10.4171/owr/2019/14}. The background describes Oberwolfach Reports as reports of mathematical workshops and their open-problem discussions. No problem-specific page or named contributor is supplied.
\item Supplied literature source, labeled ``Oberwolfach representation theory report'': \url{https://doi.org/10.4171/owr/2019/14}.
\item Literature assessment dated \textbf{2026-08-22}: the supplied curation labels the problem partially solved, reports substantial verified cases and reductions in the cited literature, and says that no general theorem covering all finite groups and the full stated \(p\)-adic assertion was found. This is supplied dated curation, not an independently verified literature assessment. The background expressly calls for independent verification of time-sensitive status and cited claims.
\end{itemize}
\end{document}
